\chapter{مثال من الواقع: متابعة مهلة تسوية الضبط التمويني}
\label{ch:settlement}

\section{القصة الحقيقية التي سنطبّقها}
\label{sec:real-story}
\textbf{لماذا؟} أتقنت في الفصل الأول بناء إجراء بموافقة وتفرّع. الآن نطبّق الأدوات نفسها على عمل يتكرر في «دائرة حماية المستهلك» كل يوم: ضبط تمويني نُظّم بحق فرن، وله مهلة تسوية، وبعدها يُحال إلى القضاء إن لم يُسوَّ.

في الواقع، عندما تجد دورية حماية المستهلك مخالفة في فرن، تنظّم «ضبطاً تموينياً» وفق قانون حماية المستهلك الصادر بالمرسوم التشريعي رقم 8 لعام 2021. لبعض المخالفات يجيز القانون التسوية أمام المديرية خلال سبعة أيام من تنظيم الضبط. فإن لم تتم التسوية خلال المهلة، أو كانت المخالفة مما لا تجوز تسويته، يُحال الضبط إلى القضاء. المشكلة العملية أن المهلة تُتابع من الذاكرة أو من دفتر، فتمر أيام بعد انتهائها قبل أن يلاحظ أحد.

\begin{table}[H]
\centering
\caption{المتابعة في الواقع وفي الإجراء الداخلي}
\label{tab:real-map}
\begin{tabularx}{\linewidth}{R{3.2cm} Y Y}
\rcHeaderRow
\thh{المرحلة} & \thh{ما يحدث في الواقع} & \thh{ما يقابله في الإجراء}\\
اعتماد الضبط & يعتمد رئيس الدائرة المحضر، فيصبح الفرن مخالفاً بانتظار التسوية. & مشغّل «تحديث سجل» مع شرط على «الحالة الرقابية».\\
بدء المهلة & يُبلَّغ موظف التسويات ورئيس الدورية بأن المهلة بدأت. & «إرسال إشعار» إلى فادي ومازن.\\
انتظار المهلة & تمر سبعة أيام. & «تأخير» بمدة سبعة أيام.\\
الفحص بعد المهلة & هل سُوّيت المخالفة خلالها؟ & «جلب سجل» ثم «تفرّع شرطي».\\
قرار الإحالة & يقرر رئيس الدائرة إحالة الضبط، أو يؤجلها لسبب مكتوب. & «طلب موافقة» بإجراءين مخصصين.\\
الإحالة & يُعَدّ ملف الإحالة ويتغير وضع الفرن. & «تحديث سجل كيان» و«إنشاء مهمة» لفادي.\\
\end{tabularx}
\end{table}

\begin{rcnote}[ما الحقيقي وما التدريبي؟]
عمل الدورية، والضبط وفق المرسوم التشريعي 8 لعام 2021، ومهلة التسوية سبعة أيام لبعض المخالفات، والإحالة إلى القضاء عند عدم التسوية: كلها مأخوذة من الواقع. أسماء الأشخاص والأفران وأرقام الضبوط تدريبية، وتحمل كلمة «تجريبي» أو «تدريب». هل تُحسب المهلة بالأيام التقويمية أم بأيام العمل؟ وأي المخالفات تجوز تسويتها؟ هذا تحدده المديرية؛ الدليل يبيّن مكان الإعداد فقط.
\end{rcnote}
\textbf{تحقق:} تستطيع أن تفرق بين ثلاثة أشياء: الضبط (ما رأته الدورية)، والمهلة (وقت يمر)، وقرار الإحالة (رأي رئيس الدائرة بعد المهلة).

\section{تأكد من الفريق والحقل الرقابي}
\label{sec:real-team}
\textbf{لماذا؟} الإجراء يكتب في حقل ويرسل إلى أشخاص. إن نقص الحقل أو أحد الخيارات، أو لم يكن للشخص تعيين نشط، يتوقف الإجراء في منتصفه.
\begin{table}[H]
\centering
\caption{الفريق في هذا المثال}
\label{tab:real-team}
\begin{tabularx}{\linewidth}{R{3.6cm} Y Y}
\rcHeaderRow
\thh{العضو ورقمه} & \thh{المنصب} & \thh{دوره في الإجراء}\\
مازن قاسم — تجريبي؛ \en{DEMO-CP-004} & رئيس دورية حماية المستهلك — 01 & يُبلَّغ ببدء المهلة لأنه نظّم الضبط.\\
سلمى عبود — تجريبي؛ \en{DEMO-CP-006} & رئيس دائرة حماية المستهلك & تقرر الإحالة أو تأجيلها بعد المهلة.\\
فادي الخوري — تجريبي؛ \en{DEMO-CP-007} & موظف التسويات والإحالة — 01 & يتابع المهلة، ويسجّل التسوية، ويُعدّ ملف الإحالة.\\
\end{tabularx}
\end{table}
\begin{enumerate}
\item افتح «الموارد البشرية» ← «التعيينات». تأكد أن للثلاثة تعيينات «نشط». إن لم يكونوا موجودين، أنشئ المناصب والأعضاء كما في الجدول، بالطريقة نفسها في الفصل الأول.
\item افتح «ادوات البناء» ← «الكيانات» ← «المنشآت التجارية». تأكد من حقل «الحالة الرقابية» من نوع «اختيار» بالخيارات الأربعة: «سليمة»، و«مخالفة — بانتظار التسوية»، و«سُوّيت المخالفة»، و«أُحيل الضبط إلى القضاء».
\item اكتب في ورقة جانبية مفتاح الحقل و«القيمة المخزّنة» لكل خيار. في الشروط أدناه نكتب \en{inspection\_status} و\en{pending\_settlement} و\en{settled} و\en{referred}؛ استبدلها بقيم مؤسستك.
\end{enumerate}
\textbf{تحقق:} يظهر الحقل في ملف «فرن الأمانة» بخياراته الأربعة، ولكل خيار قيمة مخزّنة مدوّنة عندك.

\section{اكتب القاعدة قبل البناء}
\label{sec:real-design}
\textbf{لماذا؟} هذه القاعدة أطول من قاعدة الفصل الأول، وفيها انتظار أيام. الخطأ فيها لا يظهر إلا بعد أسبوع، لذلك نصمّمها بدقة قبل البناء.

القاعدة: «عندما تتغير «الحالة الرقابية» لفرن إلى «مخالفة — بانتظار التسوية»، أبلغ فادي ومازن. انتظر سبعة أيام. ثم اقرأ ملف الفرن من جديد. إذا لم تعد الحالة «بانتظار التسوية»، فانتهِ. وإلا اطلب من سلمى قرار الإحالة. إذا قررت الإحالة، فاكتب «أُحيل الضبط إلى القضاء» وأنشئ لفادي مهمة ملف الإحالة. وإذا أجّلت، فأبلغ فادي بسببها.»

\begin{figure}[H]
\centering
\begin{tikzpicture}[node distance=0.5cm]
\node[stepBox] (a) {\RL{تحديث سجل: صار مخالفاً؟}};
\node[stepBox,below=of a] (b) {\RL{إشعار فادي ومازن}};
\node[stepBox,below=of b] (c) {\RL{تأخير سبعة أيام}};
\node[stepBox,below=of c] (d) {\RL{جلب ملف الفرن من جديد}};
\node[stepBox,below=of d] (e) {\RL{ما زال بانتظار التسوية؟}};
\node[stepBox,below=of e] (f) {\RL{قرار سلمى: إحالة أو تأجيل}};
\node[stepBox,below=of f] (g) {\RL{تحديث الحالة ومهمة فادي}};
\draw[flowsTo] (a) -- (b);
\draw[flowsTo] (b) -- (c);
\draw[flowsTo] (c) -- (d);
\draw[flowsTo] (d) -- (e);
\draw[flowsTo] (e) -- (f);
\draw[flowsTo] (f) -- (g);
\end{tikzpicture}
\rclegend
\caption{مسار متابعة المهلة}
\label{fig:real-design}
\shotnote{يُظهر المخطط الطريق الأطول فقط، حين لا تُسوّى المخالفة وتقرر سلمى الإحالة. أهم ما فيه خطوة «جلب» بعد التأخير: بيانات المشغّل عمرها سبعة أيام، والقرار يجب أن يُبنى على حال الفرن الآن.}
\end{figure}
\textbf{تحقق:} تعرف لماذا نقرأ الملف بعد التأخير بدلاً من بيانات المشغّل: لأن فادي ربما سجّل التسوية خلال المهلة.

\section{أنشئ الإجراء بمشغّل «تحديث سجل»}
\label{sec:real-trigger}
\textbf{لماذا؟} الحالة تتغير بتعديل ملف موجود، لا بإنشائه. لذلك نستخدم «تحديث سجل». لكنه يعمل مع كل تعديل، حتى تصحيح رقم هاتف، فنحتاج حارساً في أول المسار.
\begin{enumerate}
\item في مجلد «حماية المستهلك — تدريب» اضغط «إجراء جديد». الاسم «متابعة مهلة تسوية الضبط — تدريب».
\item اضغط «مزيد من المُشغِّلات» وابحث عن «تحديث سجل»، ثم «إنشاء وفتح».
\item اختر «المنشآت التجارية» في «الجدول».
\end{enumerate}
\begin{rcnote}[ما الذي يحمله المشغّل؟]
يحمل «تحديث سجل» ثلاث قيم مفيدة: \en{\$event.data} الملف بعد التعديل، و\en{\$event.oldData} الملف قبله، و\en{\$event.changedFields} الحقول التي تغيّرت. بها نعرف أن الحالة تغيّرت الآن، لا أنها كانت كذلك من قبل.
\end{rcnote}
\textbf{تحقق:} المشغّل على اللوحة يعرض «المنشآت التجارية».

\section{ضع الحارس: تابع فقط إن صار الفرن مخالفاً الآن}
\label{sec:real-guard}
\textbf{لماذا؟} بلا حارس، كل تعديل على فرن مخالف يبدأ مهلة جديدة، فيصل لفادي عشرات الإشعارات. الحارس يسمح بالمتابعة فقط عندما تنتقل الحالة إلى «بانتظار التسوية» من حالة أخرى.
\begin{table}[H]
\centering
\caption{إعداد الحارس}
\label{tab:real-guard}
\begin{tabularx}{\linewidth}{R{3.2cm} Y}
\rcHeaderRow
\thh{الإعداد} & \thh{القيمة وسببها}\\
الخطوة & \req{} «تفرّع شرطي» مباشرة بعد المشغّل. الاسم المرجعي \en{guard}.\\
المسار «بدأت المهلة» & \req{} \en{\$event.data.inspection\_status == "pending\_settlement" and \$event.oldData.inspection\_status != "pending\_settlement"}.\\
المسار البديل & \req{} «إنهاء» بالنتيجة \en{not\_relevant}.\\
\end{tabularx}
\end{table}
\begin{enumerate}
\item أضف «تفرّع شرطي» بعد المشغّل، وسمّ المسار الأول «بدأت المهلة».
\item اضغط «حدّد الشرط» ← «شرط مخصص». أدرج الحقل من «المُشغِّل» بـ«إدراج بيانات» بدلاً من كتابة المفتاح، ثم أكمل المقارنة كما في الجدول. الكلمة \en{and} تعني «و»: يجب أن يصح الجزءان.
\item على المسار البديل أضف «إنهاء».
\end{enumerate}
\textbf{تحقق:} باختبار هذه الخطوة: \en{oldData} «سليمة» و\en{data} «بانتظار التسوية» يسلك «بدأت المهلة». وإذا كانت القيمتان «بانتظار التسوية»، أو الجديدة «سُوّيت المخالفة»، يسلك البديل.

\section{أبلغ فادي ومازن ببدء المهلة}
\label{sec:real-notify}
\textbf{لماذا؟} يجب أن يعرف موظف التسويات أن العدّ بدأ اليوم، فيستقبل صاحب الفرن إن جاء للتسوية. ويعرف رئيس الدورية أن ضبطه دخل المتابعة.
\begin{enumerate}
\item على المسار «بدأت المهلة» أضف «إرسال إشعار». اختر فادي ومازن في «المستلمون».
\item «عنوان الرسالة»: «بدأت مهلة تسوية: » واسم الفرن مُدرجاً من «المُشغِّل».
\item «محتوى الرسالة»: «اعتُمد ضبط بحق هذه المنشأة. مهلة التسوية سبعة أيام من اليوم. سجّل التسوية في الملف إن تمت.»
\item الاسم المرجعي \en{notify\_start}، وفعّل «إعادة المحاولة قبل الاستسلام».
\end{enumerate}
\textbf{تحقق:} الخطوة «جاهز»، والعنوان يحمل مرجع اسم الفرن.

\section{انتظر المهلة بخطوة «تأخير»}
\label{sec:real-delay}
\textbf{لماذا؟} «تأخير» يوقف هذا التشغيل وحده مدة محددة، ثم يتابع من حيث توقف. لا يُثقل النظام، ولا يتأثر به فرن آخر؛ لكل ضبط عملية تشغيل بمهلتها.
\begin{table}[H]
\centering
\caption{إعداد «تأخير»}
\label{tab:real-delay}
\begin{tabularx}{\linewidth}{R{3.2cm} Y}
\rcHeaderRow
\thh{الحقل} & \thh{القيمة وسببها}\\
الوضع & \req{} «لمدة محددة» إذا كانت المهلة أياماً تقويمية. اختر «أيام العمل» إذا اعتمدت المديرية أيام العمل؛ تُتخطّى عندئذ العطل ونهايات الأسبوع بحسب تقويم عمل المؤسسة.\\
المدة & \req{} \en{7 days} في «لمدة محددة»، أو 7 في «أيام العمل».\\
في الساعة & \opt{} مع «أيام العمل»: 09:00، فيصل القرار في أول الدوام.\\
الاسم المرجعي & \rec{} \en{wait\_period}.\\
\end{tabularx}
\end{table}
\begin{enumerate}
\item أضف «تدفق التحكم» ← «تأخير» بعد الإشعار.
\item أدخل القيم بحسب سياسة المديرية.
\end{enumerate}
\textbf{تحقق:} تعرف أن العملية ستظهر سبعة أيام بحالة «بانتظار استجابة» عند هذه الخطوة، وأن هذا ليس خطأً.

\section{اقرأ ملف الفرن من جديد بعد المهلة}
\label{sec:real-get}
\textbf{لماذا؟} \en{\$event.data} صورة الملف لحظة اعتماد الضبط. إن سجّل فادي التسوية في اليوم الثالث، فالصورة القديمة لا تعرف ذلك. «جلب سجل» يقرأ الملف كما هو الآن.
\begin{enumerate}
\item أضف «السجلات» ← «جلب سجل» بعد التأخير.
\item «الجدول» «المنشآت التجارية»، و«السجل» بتعبير \en{\$event.data.id}.
\item «إذا كان معرّف السجل فارغاً»: «أفشِل الخطوة».
\item الاسم المرجعي \en{recheck}. افتح «متقدم» ← «اتباع مسار الفشل»، وعلى مسار الفشل أضف إشعاراً لسلمى: «تعذّرت قراءة ملف منشأة بعد مهلة التسوية؛ ربما حُذف. راجع الضبط يدوياً.»
\end{enumerate}
\textbf{تحقق:} يخرج من \en{recheck} خطان. الحذف بين الضبط ونهاية المهلة لن يُسقط المتابعة بصمت.

\section{تفرّع: هل ما زال الفرن بانتظار التسوية؟}
\label{sec:real-branch}
\textbf{لماذا؟} الطلب يُرسل إلى سلمى فقط إن لم تتم التسوية. أما إن سُوّيت، فلا عمل لأحد، وينتهي التشغيل بنتيجة واضحة.
\begin{enumerate}
\item أضف «تفرّع شرطي» بعد \en{recheck}، والاسم المرجعي \en{still\_pending}.
\item المسار «لم تُسوَّ»: اضغط «حدّد الشرط» ← «شرط مخصص»، ثم «إدراج بيانات» ← خطوة \en{recheck} ← حقل «الحالة الرقابية»، وأكمل بـ\en{== "pending\_settlement"}.
\item المسار البديل: «إنهاء» بالنتيجة \en{settled\_or\_changed}.
\end{enumerate}
\begin{rcnote}[إن لم تظهر حقول \en{recheck}]
قائمة «إدراج بيانات» تعرض حقول الخطوة بعد أن تعمل في اختبار. شغّل اختبار خطوة \en{recheck} على فرن تدريبي، ثم عد إلى الشرط.
\end{rcnote}
\textbf{تحقق:} باختبار التفرّع: حالة «بانتظار التسوية» تسلك «لم تُسوَّ»، وحالة «سُوّيت المخالفة» تسلك البديل.

\section{اطلب قرار سلمى بإجراءات مخصّصة}
\label{sec:real-approval}
\textbf{لماذا؟} الإجراءان الجاهزان، الموافقة وطلب التغييرات، لا يطابقان لغة الدائرة. سلمى تقرر «إحالة إلى القضاء» أو «تأجيل الإحالة». نعدّل «إجراءات الطلب» لتحمل الأزرار هذين الاسمين، ولكل منهما نتيجة يقرؤها التفرّع.
\begin{table}[H]
\centering
\caption{إجراءات الطلب المخصّصة}
\label{tab:real-actions}
\begin{tabularx}{\linewidth}{R{3cm} Y Y}
\rcHeaderRow
\thh{الإعداد} & \thh{الإحالة} & \thh{التأجيل}\\
المفتاح & \en{refer} & \en{postpone}\\
التسمية & «إحالة إلى القضاء» & «تأجيل الإحالة»\\
النتيجة & \en{referred} & \en{postponed}\\
التقدم & «إكمال» & «إكمال»\\
التوقيع & «مطلوب»؛ قرار له أثر قانوني. & «لا شيء».\\
التعليق & «اختياري». & «مطلوب»؛ يجب أن يُكتب سبب التأجيل.\\
\end{tabularx}
\end{table}
\begin{table}[H]
\centering
\caption{بقية حقول «طلب موافقة»}
\label{tab:real-approval}
\begin{tabularx}{\linewidth}{R{4cm} Y}
\rcHeaderRow
\thh{الحقل} & \thh{القيمة}\\
عنوان الموافقة & \req{} «انتهت مهلة التسوية: » واسم الفرن من \en{recheck}.\\
تفاصيل الموافقة & \rec{} «انقضت سبعة أيام على اعتماد الضبط ولم تُسجَّل تسوية. قرّري الإحالة إلى القضاء أو تأجيلها مع ذكر السبب.»\\
مراحل الموافقة & \req{} مرحلة واحدة؛ سلمى معتمدة.\\
المهلة الزمنية & \rec{} \en{3 days}.\\
إرفاق أحدث إصدار من السجل & \rec{} مفعّل، والسجل \en{\$event.data.id}؛ ترى سلمى الملف الذي تقرر عليه.\\
الاسم المرجعي & \req{} \en{referral}.\\
\end{tabularx}
\end{table}
\begin{enumerate}
\item على المسار «لم تُسوَّ» أضف «طلب موافقة».
\item في «إجراءات الطلب» عدّل الإجراء الأول بقيم عمود الإحالة، والثاني بقيم عمود التأجيل. لا تترك تسمية أو نتيجة فارغة؛ يرفضها المحرر بـ«مفتاح النتيجة مطلوب.»
\item أدخل بقية الحقول من الجدول الثاني.
\item في «مرفقات المكتبة» اضغط زر التعبير واكتب \en{[]}، كما في الخطوة \ref{sec:attachments} من الفصل الأول. بدونها يفشل الطلب برسالة «Internal command failed with status 422».
\end{enumerate}
\textbf{تحقق:} في المراجعة لا مشكلات على الخطوة، وأزرار الطلب ستحمل الاسمين الجديدين، وتحت «مرفقات المكتبة» القيمة \en{[]}.

\section{نفّذ القرار: إحالة أو تأجيل}
\label{sec:real-outcome}
\textbf{لماذا؟} قرار سلمى يجب أن يظهر في ملف الفرن وفي عمل فادي، وإلا بقي داخل الطلب.
\begin{enumerate}
\item أضف «تفرّع شرطي» بعد \en{referral} بثلاثة مسارات: «إحالة» بشرط \en{\$steps.referral.output.terminalOutcomeKey == "referred"}، و«تأجيل» بشرط \en{\$steps.referral.output.terminalOutcomeKey == "postponed"}، والبديل.
\item على «إحالة»: «تحديث سجل كيان» للسجل \en{\$event.data.id}، و«الحالة الرقابية» = «أُحيل الضبط إلى القضاء». ثم «إنشاء مهمة» لفادي بعنوان «إعداد ملف إحالة: » واسم الفرن، وموعد بعد يومي عمل.
\item على «تأجيل»: «إرسال إشعار» لفادي: «أجّلت رئيسة الدائرة الإحالة. افتح الطلب لقراءة السبب.» ثم «إنهاء» بالنتيجة \en{postponed}.
\item على البديل، وهو انتهاء مهلة سلمى دون رد: إشعار لسلمى وفادي: «انتهت مهلة قرار الإحالة دون رد.» ثم «إنهاء».
\end{enumerate}
\begin{rcwarn}[لا تجعل الإجراء يحيل وحده]
قد يبدو أسهل أن يكتب الإجراء «أُحيل الضبط إلى القضاء» بعد سبعة أيام مباشرة. لا تفعل: قد تكون التسوية تمت ولم تُسجَّل، أو تكون المخالفة مما يُقدَّر. الإجراء يذكّر ويجهّز، والقرار لصاحب الصلاحية.
\end{rcwarn}
\textbf{تحقق:} قائمة «الخطوات» «كل الخطوات جاهزة»، ولكل مسار نهاية.

\section{اختبر المسارات الأربعة قبل النشر}
\label{sec:real-test}
\textbf{لماذا؟} هذا الإجراء ينتظر أسبوعاً قبل أن يكشف خطأه. اختبر كل مسار الآن بالمحاكاة واختبار الخطوات، فلا تنتظر أسبوعاً لتكتشف شرطاً خاطئاً.
\begin{table}[H]
\centering
\caption{حالات الاختبار ونتائجها المتوقعة}
\label{tab:real-tests}
\begin{tabularx}{\linewidth}{R{4.4cm} Y}
\rcHeaderRow
\thh{البيانات} & \thh{النتيجة المتوقعة}\\
تعديل هاتف فرن مخالف؛ الحالة قبل وبعد «بانتظار التسوية» & ينتهي عند الحارس بـ\en{not\_relevant}. لا إشعار.\\
تغيير الحالة من «سليمة» إلى «بانتظار التسوية» & إشعار، ثم يقف عند «تأخير».\\
\en{recheck} يعيد «سُوّيت المخالفة» & ينتهي بـ\en{settled\_or\_changed}. لا طلب لسلمى.\\
\en{recheck} يعيد «بانتظار التسوية»، ونتيجة \en{referral} \en{referred} & تحديث الحالة ومهمة فادي.\\
نتيجة \en{referral} \en{postponed} & إشعار فادي، ثم \en{postponed}.\\
\end{tabularx}
\end{table}
\begin{enumerate}
\item في «بيانات الاختبار» اكتب كائناً فيه \en{data} و\en{oldData} لكل حالة من الحالتين الأوليين، وشغّل «محاكاة (آمنة)».
\item للحالات الثلاث الأخيرة اختبر التفرّعين \en{still\_pending} والتفرّع بعد \en{referral} بـ«اختبار هذه الخطوة»، بقيم مكتوبة لمخرجات الخطوات السابقة.
\end{enumerate}
\textbf{تحقق:} طابقت كل حالة في الجدول نتيجتها. إن سلك تعديل الهاتف طريق المهلة، فالحارس لا يقارن \en{oldData}.

\section{انشر وطبّق على «فرن الأمانة»}
\label{sec:real-live}
\textbf{لماذا؟} الآن نشاهد الإجراء مع عمل حقيقي: دورية تنظّم ضبطاً، ورئيسة الدائرة تعتمده، فيبدأ العدّ وحده.
\begin{enumerate}
\item اضغط «مراجعة ونشر» ← «نشر \en{v1}».
\item اعتمد بحساب سلمى محضراً تدريبياً فيه مخالفة على «فرن الأمانة»، أو غيّر «الحالة الرقابية» يدوياً من «سليمة» إلى «مخالفة — بانتظار التسوية» في ملف الفرن.
\item بحساب فادي افتح الإشعارات: «بدأت مهلة تسوية: فرن الأمانة».
\item بحساب المسؤول افتح «عمليات التشغيل» ← «فعلي».
\end{enumerate}
\begin{rcnote}[للتدريب: لا تنتظر أسبوعاً]
لتشاهد بقية المسار اليوم، انسخ الإجراء في مسودة تدريبية أخرى واجعل «المدة» \en{5 minutes}، وانشرها وحدها، ثم أعد المدة وأوقف النسخة بعد التدريب. لا تغيّر مدة الإجراء الذي يتابع ضبوطاً حقيقية.
\end{rcnote}
\textbf{تحقق:} عملية واحدة «بانتظار استجابة» عند «تأخير»، فيها موعد الاستئناف. وعدّل هاتف الفرن: لا تبدأ عملية ثانية.

\section{سجّل تسوية خلال المهلة وراقب الإجراء ينسحب}
\label{sec:real-settle}
\textbf{لماذا؟} أهم ما يثبت صحة القاعدة أن تبقى صامتة حين يكون العمل قد تم.
\begin{enumerate}
\item بحساب فادي افتح ملف الفرن، وغيّر «الحالة الرقابية» إلى «سُوّيت المخالفة»، ثم احفظ.
\item بعد انتهاء المدة، افتح العملية في «عمليات التشغيل».
\end{enumerate}
\textbf{تحقق:} قرأت \en{recheck} «سُوّيت المخالفة»، وانتهت العملية بـ\en{settled\_or\_changed}، ولم يصل سلمى أي طلب.

\section{جرّب طريق الإحالة حتى نهايته}
\label{sec:real-refer}
\textbf{لماذا؟} طريق الإحالة له أثر على ملف الفرن وعمل فادي، فيجب أن يُرى مرة حقيقية قبل الاعتماد عليه.
\begin{enumerate}
\item أنشئ ضبطاً تدريبياً على «فرن الجامعة — تدريب» ودعه بلا تسوية حتى تنتهي المدة التدريبية.
\item بحساب سلمى افتح «طلبات بانتظارك» ← «انتهت مهلة التسوية: فرن الجامعة — تدريب». الزران «إحالة إلى القضاء» و«تأجيل الإحالة».
\item اضغط «إحالة إلى القضاء» ووقّع.
\item بحساب فادي افتح مهامه وملف الفرن.
\end{enumerate}
\textbf{تحقق:} «الحالة الرقابية» «أُحيل الضبط إلى القضاء»، ولفادي مهمة «إعداد ملف إحالة: فرن الجامعة — تدريب»، والعملية «مكتمل». كرر مع «تأجيل الإحالة»: التعليق إلزامي، ويصل فادي الإشعار، ولا تتغير الحالة.

\section{تابع الإجراء أسبوعياً كمسؤول}
\label{sec:real-monitor}
\textbf{لماذا؟} إجراء يعمل أسابيع يحتاج نظرة دورية: كم ضبطاً في المهلة الآن؟ هل فشل شيء؟ هل تأخرت سلمى في قرار؟
\begin{enumerate}
\item افتح «الإجراءات الداخلية» واختر «تحتاج إلى متابعة». إن ظهر الإجراء، فاضغط «فتح التشغيل الفاشل».
\item افتح الإجراء ← «عمليات التشغيل» ← «فعلي». العمليات عند «تأخير» ضبوط في المهلة؛ والعمليات عند \en{referral} قرارات تنتظر سلمى.
\item إن تغيّرت سياسة المهلة، عدّل «تأخير» وانشر إصداراً جديداً. الضبوط التي بدأت مهلتها تكمل على إصدارها القديم، وهذا صحيح قانونياً لأن مهلتها حُسبت عند بدئها.
\end{enumerate}
\textbf{تحقق:} تستطيع أن تقول من «عمليات التشغيل» وحدها: كم فرناً في المهلة، وكم قراراً معلقاً، ومتى تنتهي أقرب مهلة.

\section{قائمة البطل: ما الذي يجب أن تستطيعه الآن}
\label{sec:real-hero}
\textbf{لماذا؟} تصبح مرجع فريقك في الإجراءات عندما تشرح لماذا سلك كل تشغيل طريقه. أجب عن هذه الأسئلة من التطبيق، دون الرجوع إلى الدليل.
\begin{table}[H]
\centering
\caption{أثبت إتقانك}
\label{tab:real-hero}
\begin{tabularx}{\linewidth}{Y Y}
\rcHeaderRow
\thh{السؤال} & \thh{أين تجد الجواب؟}\\
لماذا لم يبدأ تعديل الهاتف مهلة جديدة؟ & مسار الحارس البديل في العملية، ونتيجتها \en{not\_relevant}.\\
متى تنتهي مهلة «فرن الأمانة»؟ & موعد الاستئناف في خطوة «تأخير» للعملية المنتظرة.\\
لماذا لم يصل سلمى طلب عن فرن سُوّي؟ & مخرجات \en{recheck}، والمسار البديل للتفرّع.\\
من قرر إحالة «فرن الجامعة» ومتى؟ & ردود خطوة \en{referral}، والطلب في «مكتملة» عند سلمى.\\
لماذا لم يتأثر ضبط قديم بتغيير المدة؟ & رقم الإصدار في العملية، و«سجل الإصدارات».\\
\end{tabularx}
\end{table}
\begin{itemize}
\item أضف «إرسال تذكير بالموافقة» لسلمى إذا لم ترد خلال يوم، ثم اختبره.
\item ابنِ إجراءً مجدولاً يرسل لسلمى كل أحد عدد الأفران «بانتظار التسوية» بـ«البحث عن سجلات كيان»، كما في الفصل الأول.
\item اشرح لزميل الفرق بين بيانات المشغّل ونتيجة «جلب سجل»، من مثال التسوية خلال المهلة.
\end{itemize}
\textbf{تحقق:} أجبت عن الأسئلة الخمسة من التطبيق، ونفّذت تمريناً واحداً على الأقل دون الدليل.
